\section{A filtered cut--join identity}\label{sec:cutjoin}

We now construct a degree-lowering operation which removes a long part of
a loop.  Its essential property is an interchange identity with the loop
product, with an additive length cost independent of the degree.  Applied to the
sphere's product powers, this identity makes the cut of a suitable cycle
equal to the constant-loop fundamental class.  These constructions apply
to every smooth metric on $S^n$, $n\ge3$; no assumption about the number
of closed geodesics is made in this section.

We retain $\mathcal A_a=\{\ell<a\}$ and $d=n-1$.  A chain is always a finite
singular chain with coefficients in $\kk$.  Writing $c\simeq c'$ in a
subspace means that their difference bounds a chain in that subspace.

\subsection{The sphere classes and their length thresholds}

Write $\bullet$ for the Chas--Sullivan homology product
\cite{ChasSullivan1999}, of degree $-n$,
and $\circledast$ for the Goresky--Hingston product on
$H^*(\Lambda,M)$ \cite{GoreskyHingston2009}, of degree $n-1$.
We use the following precise
consequences of sphere resonance and the calculation of these products.

\begin{proposition}[Sphere powers and resonance]\label{prop:sphere-powers}
There are constants $\alpha>0$ and $C_H$ depending on $g$ such that
\begin{equation}\label{eq:cj-resonance}
 \bigl|k-\alpha\operatorname{cr}(z)\bigr|\le C_H
\end{equation}
for every nonzero homogeneous ordinary homology or cohomology class $z$
of degree $k$.  For homology, $\operatorname{cr}$ is the infimum of the
levels $a$ carrying the class in $X_a=\{F<a\}$; for cohomology it is
the supremum of the levels on which its restriction vanishes.
Moreover, each ordinary homology and cohomology group in a fixed degree
has dimension at most one.  There are classes
\[
 \Theta\in H_{3n-2}(\Lambda),\qquad
 \omega\in H^{n-1}(\Lambda,M),\qquad
 \theta\in H^{3n-2}(\Lambda,M)
\]
with $\langle\theta,\Theta\rangle=1$ such that, for every $m\ge1$,
\begin{align}
 |\Theta^{\bullet m}|&=n+2md,&
 |\omega^{\circledast m}|&=(2m-1)d,\label{eq:cj-power-degrees}\\
 \bigl\langle\theta\circledast\omega^{\circledast m},
          \Theta^{\bullet(m+1)}\bigr\rangle&=1.
 \label{eq:cj-power-pairing}
\end{align}
In particular these powers are nonzero, and
$H_n(\Lambda)=\kk\,[M]$.
\end{proposition}

\begin{proof}
The resonance statement is Theorem~1.1 of
\cite{HingstonRademacher2013}.  Its cohomological critical levels
may be computed in singular cohomology: the open energy sublevels
and $\Lambda$ are paracompact, locally contractible Hilbert
manifolds, where singular and \v{C}ech cohomology agree.
Their Lemma~5.3 gives the rank bound;
Lemma~5.4 and Equation~(12) give the powers and the pairing, with the
index in that equation replaced by $m+1$.  These statements apply for
every $n>2$ and every field.  The indicated powers and pairing may also
be taken integrally and reduced modulo two: the integral pairing is
one, so its reduction is one in both parities of $n$.
The absolute class of degree $3n-2$ has a unique relative lift because
$3n-2>n+1$.  Finally evaluation retracts the inclusion of constant
loops, so $[M]\ne0$; the rank bound gives the last assertion.
\end{proof}

For sufficiently large $m$, set
\begin{equation}\label{eq:cj-thresholds}
 q=2md,\qquad p=q-d,\qquad
 l=\frac{n+q+C_H}{\alpha}+2,\qquad
 L=\frac{p-C_H}{\alpha}-2>0.
\end{equation}
We also require $p>n+1$.  Choose a cycle $S$ representing
$\Theta^{\bullet m}$ in $\mathcal A_l$.  Such a cycle exists already in
$X_l$ by \eqref{eq:cj-resonance}.  Choose a normalized cocycle $T$
representing $\omega^{\circledast m}$ and vanishing on $\mathcal A_L$.
Indeed its absolute class, identified with the relative class since
$p>n+1$, vanishes on $X_L$ by resonance.  The cohomology exact sequence
gives a cocycle relative to $X_L$.  Pull it back by the arclength
map of Lemma~\ref{lem:arclength}; this sends $\mathcal A_L$ into $X_L$, fixes
constants, and is homotopic to the identity.  Normalization preserves
the relative condition.  Notice the fixed difference
\begin{equation}\label{eq:cj-fixed-difference}
 l-L=\frac{2n-1+2C_H}{\alpha}+4,
\end{equation}
independent of $m$.  This is the length available for the eventual
degree reduction.

\subsection{Supported caps on small chains}

We use unnormalized singular chains and normalized cochains.  The
shuffle cross product is associative and, over $\kk$, symmetric under
permutation of factors.  Our cap evaluates a cochain on the front
face and retains the back face.  Thus successive caps correspond to
cup products in the same order, and
\[
 \partial(c\frown b)=(\partial c)\frown b+c\frown\delta b.
\]
The degenerate chains form an acyclic subcomplex in every subspace;
the usual normalization homotopies preserve simplex images.
Consequently an equality modulo degenerate chains between cycles gives
a homology in the same length sublevel.

All tubes, arcs and rotations in the construction use a fixed round
metric on $M$; all length estimates use $g$.  Put
\[
 F_\eta=\{(z,z')\in M^2:d_{\rm rd}(z,z')\ge\eta\}.
\]
A cochain relative to $F_\eta$ vanishes on every simplex contained
in this set.  Its nonzero evaluation on a face therefore supplies a
point of that face with distance less than $\eta$.  Subdividing a
finite parameter family sufficiently finely makes the corresponding
distance less than $2\eta$ on the entire parent simplex.  For product
chains we impose this condition on the whole product rectangles
before shuffling.  Maps used below are defined on the larger
$4\eta$ tubes, so uniform continuity on an intermediate closed tube
allows successive support conditions to be imposed together.

The following observation also deals with cochains whose maps are
defined only on a tube.

\begin{lemma}[Supported cochain calculus]\label{lem:cj-small-chains}
Let $U$ be a pulled-back Thom cocycle, and let $h$ be a continuous
map from its larger tube $O$ to a space carrying a cocycle $T$.
Extend $h^*T$ by zero on simplices not wholly contained in $O$, and
write $b=U\smile\widetilde{h^*T}$.  Subdivide the finite parameter
families factorwise so that a nonzero $U$-face forces its entire
product rectangle into a fixed closed intermediate tube contained
in $O$.  Its preimage in each finite compact parameter domain is
compact.  Then $\delta b=0$ on the shuffled simplices and every
later face or image-preserving comparison prism in these rectangles.

If $T$ vanishes on loops of length less than $L$, the subdivision
may also ensure the following: a contributing term with a nonzero
$U$-face and a nonzero $T$-face has $\ell(h)\ge L-\varepsilon$
throughout its whole rectangle, for any fixed $\varepsilon>0$.
The same conclusion holds for the block and skip faces used by
cup-one, and for all later image-preserving subdivision prisms.
\end{lemma}

\begin{proof}
Extension by zero defines a cochain, although generally not a
cocycle.  The ordinary cochain identity nevertheless gives
\[
 \delta b=U\smile\delta\widetilde{h^*T}.
\]
If its front $U$-face is zero, the evaluation is zero.  Otherwise the
entire parent lies in $O$, and every face entering the remaining
coboundary is an actual pullback of $T$.  They cancel because
$\delta T=0$.  This proves the assertion without discarding any
extension-by-zero terms separately.  Each block or skip face is
contained in its parent image and has the same property.

In a contributing term the $U$-face first puts the whole rectangle
in the closed intermediate tube, where $h$ is defined and uniformly
continuous on the compact parameter set.  Its nonzero $T$-face
cannot have its entire image in $\{\ell<L\}$, and hence contains
a point with $\ell(h)\ge L$.  Refining the factors so that
$\ell\circ h$ oscillates by less than $\varepsilon$ on each such
rectangle proves the estimate.  The tube conditions follow in the
same way from the gap between support and intermediate radii.
Include all homotopy intervals among the factors before refining.
The finitely many conditions then hold simultaneously.  All later
subdivision prisms and block or skip faces stay in the original
rectangles, so both the domain and length conditions persist.
\end{proof}

We will always use such small chains.  Boundary comparisons have the
following concrete meaning.  Subdivide a finite homology adequately,
then apply the cap boundary formula.  Its subdivided boundary is
compared with any earlier adequate representative by the subdivision
prism inside the original simplex images.  Factor subdivisions are
compared by crossing these prisms with the other factors.  For a
homotopy of evaluation or output maps, cross with its interval
\emph{before} capping and refine all factors, so the endpoint chains
are the original shuffled constructions.  This proves compatibility
with refinements and finite fillings at the same whole-parent length
bounds.  A single subdivision depth on the entire loop space is
neither chosen nor needed.
The shuffle step counts below are applied to these factorwise
refined product chains.  Comparison prisms need not have that
form; they are handled by the supported boundary formula instead.

\subsection{The cut and the join}

Fix $\varepsilon=10^{-3}$.  Choose a small cut radius $\eta_c$ so
that a short round arc at distance less than $4\eta_c$ has
$g$-length less than $\varepsilon$.  The diagonal Thom class has
a normalized representative
\begin{equation}\label{eq:cj-diagonal-cocycle}
 U_c\in C^n(M^2,F_{\eta_c}\cup\Delta M).
\end{equation}
To see this, its restriction to $\Delta M$ is the mod-two Euler
class of $TM$, which vanishes because $\chi(S^n)$ is even.  The
exact sequence of the triple
$F_{\eta_c}\subset F_{\eta_c}\cup\Delta M\subset M^2$
therefore lifts the Thom class to this relative group.  A relative
cocycle representative, followed by normalization, gives
\eqref{eq:cj-diagonal-cocycle}.

For paths $\gamma,\zeta$ with matching endpoints, let
$\gamma\#\zeta$ denote concatenation with half the time assigned
to each.  Brackets $[z,z']$ denote a short round arc.  On the cut
tube define
\begin{equation}\label{eq:cj-cut-maps}
 \begin{split}
 C_1(x,t)&=x[0,t]\#[x(t),x(0)],\\
 C_2(x,t)&=[x(0),x(t)]\#x[t,1],
 \end{split}
\end{equation}
where restrictions of a path are rescaled to the unit interval.
Figure~\ref{fig:native-cut} shows the tested prefix and retained
complement, with the short arcs which close them.
These maps are jointly continuous in $H^1$, including $t=0,1$.
Indeed the derivative energy of a restriction to $[s,t]$ rescaled
to unit time is
$(t-s)\int_s^t|x'|^2$, which tends to zero at a collapsing
interval, uniformly along an $H^1$-convergent family.

\begin{figure}[tbp]
\centering
\begin{tikzpicture}[x=.94cm,y=.94cm,>=Stealth,
  every node/.style={font=\small},
  first arc/.style={blue!65!black,line width=1pt},
  second arc/.style={teal!70!black,line width=1pt},
  stick/.style={gray,dashed,line width=.8pt}]
  \draw[first arc,->] (0,0) .. controls (.3,1.45) and (2.5,1.45) .. (2.8,0);
  \draw[second arc,->] (2.8,0) .. controls (2.5,-1.45) and (.3,-1.45) .. (0,0);
  \fill (0,0) circle (1.5pt);
  \fill (2.8,0) circle (1.5pt);
  \node[left=4pt] at (0,0) {$z$};
  \node[right=4pt] at (2.8,0) {$y$};
  \node[above=3pt,text=blue!65!black] at (1.4,1.08) {$x[0,t]$};
  \node[below=3pt,text=teal!70!black] at (1.4,-1.08) {$x[t,1]$};
  \draw[->] (3.5,0) -- (4.3,0);
  \draw[first arc,->] (5,0) .. controls (5.3,1.45) and (7.5,1.45) .. (7.8,0);
  \draw[stick,->] (7.8,-.08) -- (5,-.08);
  \fill (5,0) circle (1.5pt);
  \fill (7.8,0) circle (1.5pt);
  \node[left=3pt] at (5,0) {$z$};
  \node[right=3pt] at (7.8,0) {$y$};
  \node at (6.4,-.58) {$C_1(x,t)$};
  \node at (6.4,-1.05) {tested by $T$};
  \node at (8.85,0) {and};
  \draw[stick,->] (10,.08) -- (12.8,.08);
  \draw[second arc,->] (12.8,0) .. controls (12.5,-1.45) and (10.3,-1.45) .. (10,0);
  \fill (10,0) circle (1.5pt);
  \fill (12.8,0) circle (1.5pt);
  \node[left=3pt] at (10,0) {$z$};
  \node[right=3pt] at (12.8,0) {$y$};
  \node at (11.4,.55) {$C_2(x,t)$};
  \node at (11.4,1.02) {retained};
\end{tikzpicture}
\caption{The native cut at $y=x(t)$, with $z=x(0)$. Short round arcs
(dashed) close the prefix and the remaining path. The cut tests the first
loop and retains the second. The drawing separates the nearby endpoints
to show the two arcs.}
\label{fig:native-cut}
\end{figure}


For a normalized $p$-cocycle $T$ vanishing on $\mathcal A_L$, put
$q=n+p-1$ and define on adequate chains
\begin{equation}\label{eq:cj-cut-definition}
 D_Tx=(C_2)_\#\bigl((x\times I)\frown
 U_c(x(0),x(t))\frown T(C_1(x,t))\bigr).
\end{equation}
Pullback notation in this formula is abbreviated; the second cap
is defined on the supported cut tube.
Absolute extensions of the loop coproduct and their Euler-characteristic
boundary correction were studied by Cieliebak, Hingston and
Oancea~\cite[Section~8.1]{CieliebakHingstonOancea2023}.
Their odd-sphere calculation also exhibits constant-loop output terms
\cite[Proposition~8.3 and Remark~8.4]{CieliebakHingstonOancea2023}.
The filtered comparison for the cut defined here is proved below.

For joining, choose an arbitrarily small radius $\eta_j$ and a
normalized diagonal Thom cocycle $U_j$ relative to $F_{\eta_j}$.
For close $z,z'$, let $J_{z,z'}$ be the shortest plane rotation
taking $z$ to $z'$, and the identity at equality.  The rotations
$J_{z,z'}^u$, $-1\le u\le1$, depend smoothly on the two points
and converge smoothly to the identity as the points approach.
Their $g$-length distortion is at most $1+\kappa$, with
$\kappa\to0$ as $\eta_j\to0$.  Define
\begin{equation}\label{eq:cj-join-definition}
 \mu(x,y)=\bigl((J_{x(0),y(0)}x)\#y\bigr)_\#
 \bigl((x\times y)\frown U_j(x(0),y(0))\bigr).
\end{equation}

\begin{lemma}[Boundary and length properties]\label{lem:cj-operations}
The cut has degree $-q$, takes absolute cycles to absolute cycles,
and preserves global null-homology.  For a cycle in $\mathcal A_a$ it is
supported in
\begin{equation}\label{eq:cj-cut-bound}
 \mathcal A_{a-L+1}.
\end{equation}
The same cut estimate holds on finite fillings.  The join has degree
$-n$ and satisfies the product boundary formula.  On any prescribed
finite collection of chains and fillings its radius can be chosen
so that its output length is less than the sum of the two input
length bounds plus $1$.  All these statements are compatible with
the refinement comparisons above.

If $c$ is a fundamental cycle of constant loops, then
\begin{equation}\label{eq:cj-chain-unit}
 \mu(c,x)\simeq x\quad\text{in }\mathcal A_a
\end{equation}
for every cycle $x$ in $\mathcal A_a$.
\end{lemma}

\begin{proof}
The cut degree is $1-n-p=-q$.  The cap boundary formula applies
by Lemma~\ref{lem:cj-small-chains}; its time-end terms vanish
because $(x(0),x(t))$ is exactly diagonal at $t=0,1$.  Applying
the same construction to a finite filling proves null-homology
preservation.  On a surviving parent the tested first loop has
length at least $L-\varepsilon$, while
\[
 \ell(C_1(x,t))+\ell(C_2(x,t))
 \le\ell(x)+2\varepsilon.
\]
This gives \eqref{eq:cj-cut-bound}.  If its upper level is
nonpositive, no term survives.  The join boundary formula uses
the closed cocycle $U_j$.  Its output length is at most
$(1+\kappa)\ell(x)+\ell(y)$, and $\eta_j$ may be chosen
after all finite length bounds have been specified.

For the unit, consider a shuffle of $c\times x$.  A constant-factor
step remaining after the $n$-cap makes the output degenerate:
the rotated constant is exactly $x(0)$, so the output is independent
of that factor.  The remaining terms have all $n$ constant-factor
steps first.  Their coefficient is the Thom evaluation on
$[M]\times\{x(0)\}$, at the initial parameter vertex of the
$x$-simplex, and equals one.  Thus the output is a constant half
followed by $x$.  Removing the waiting half by reparametrization
gives \eqref{eq:cj-chain-unit} without changing length.
\end{proof}

\subsection{Interchanging which factor is cut}

\begin{theorem}[Filtered interchange]\label{thm:filtered-interchange}
Let $T$ be a normalized cocycle of degree $p\ge1$ vanishing on
$\mathcal A_L$, where $L>0$, and put $q=n+p-1$.  Let $x$ and $y$ be
arbitrary absolute cycles of respective degrees $j,N$ with
\[
 j>n+p,\qquad N>n+p,
\]
supported in $\mathcal A_a$ and $\mathcal A_b$.  With sufficiently small join radius
and adequate refinements,
\begin{equation}\label{eq:cj-interchange}
 \mu(D_Tx,y)\simeq\mu(D_Ty,x)
 \quad\text{in }\mathcal A_{a+b-L+c_{\rm cj}},\qquad c_{\rm cj}=4.
\end{equation}
Both sides have degree $j+N-q-n$.  The constant $c_{\rm cj}$
is independent of both degrees, the cocycle, and all length bounds.
The join radius may be chosen for any additional prescribed finite
collection of fillings without changing this constant.
\end{theorem}

\begin{proof}
Write $J=J_{x(0),y(0)}$ and $Z=(Jx)\#y$.  On
$x\times I_s\times y$, first cap with
$\sigma=U_j(x(0),y(0))$, then with
$U_c(Z(0),Z(s))$, then with $T(C_1(Z,s))$, and output
$C_2(Z,s)$.  Call this cycle $P(x,y)$.  Its time endpoints
vanish by \eqref{eq:cj-diagonal-cocycle}.  We prove
\begin{equation}\label{eq:cj-three-comparisons}
 \mu(D_Tx,y)\simeq P(x,y)\simeq P(y,x)
                   \simeq\mu(D_Ty,x).
\end{equation}
All homologies will have the bound in \eqref{eq:cj-interchange}.

\paragraph{The second half vanishes.}
Split $I_s$ at $1/2$.  At the seam $Z(1/2)=Z(0)$, so each
piece is separately a cycle after capping.  For $s\ge1/2$, both
the cut evaluation and the second-cut output depend only on
$(y,s)$.  A shuffle contains $j$ steps from $x$.  Its join face
uses at most $n$ of them, and its $T$-face at most $p$.
Since $j>n+p$, another such step belongs either to the cut Thom
face or to the retained output face.  In the first case the
normalized Thom evaluation is zero; in the second the output
simplex is degenerate.  In fact projection drops the unused
$x$-step, and the output factors through the corresponding
standard degeneracy map.  Deleting a repeated projected vertex
uses a face of the original simplex, so this factorization stays
within the cut domain.  The second-half cycle therefore bounds
in its image sublevel.

\paragraph{The first half gives the native cut.}
Put $t=2s$ on the first half.  The first cut is exactly
$JC_1(x,t)$.  The second cut consists of an arc, the rotated
$x$-tail and $y$, with respective durations
\[
 \frac12,\qquad\frac{1-t}{2(2-t)},\qquad\frac1{2(2-t)}.
\]
Interpolate them linearly to $1/4,1/4,1/2$.  The result is
$(JC_2(x,t))\#y$.  Length is unchanged, and continuity at
$t=1$ follows because the shrinking tail contributes at most
a fixed multiple of $\int_t^1|x'|^2$ to the energy.

Keep this output fixed and move the cut evaluation and tested
first loop from those of $Jx$ to those of $J^u x$, with $u$
decreasing from $1$ to $0$.  Cross with this parameter before
capping.  Round rotations preserve the cut distance, so all
surviving parents support the fixed output throughout.  The cap
boundary formula gives a homology: it requires continuous
supported maps, not an identity between the fixed output and an
intermediate tested cut.  Its time endpoints remain diagonal.
Small rotation distortion keeps the first-cut plus output length
within a fixed small error of $\ell(x)+\ell(y)$.

At $u=0$ put
\[
 \gamma=U_c(x(0),x(t))\smile T(C_1(x,t)),
\]
using extension by zero as in Lemma~\ref{lem:cj-small-chains}.
This normalized cochain has degree $n+p$ and is independent of
$y$.  Impose the native cut smallness even before the join cap.
We may then commute $\sigma$ and $\gamma$ using
\begin{equation}\label{eq:cj-cup-one}
 \delta(\sigma\smile_1\gamma)
 =\sigma\smile\gamma+\gamma\smile\sigma
       +\sigma\smile_1\delta\gamma.
\end{equation}
For completeness, if the degrees of cochains $u,v$ are $r,s$,
their mod-two cup-one on a simplex of dimension $r+s-1$ is
\[
 \sum_{i=0}^{r-1}
 u[0,\ldots,i,i+s,\ldots,r+s-1]\,
 v[i,\ldots,i+s].
\]
The coboundary formula follows by pairing deletions outside and
inside the block; the shared-end deletions cancel in adjacent
summands except for the two outer cups.  This gives
\eqref{eq:cj-cup-one} since $\delta\sigma=0$.

Lemma~\ref{lem:cj-small-chains} makes $\delta\gamma$ zero
on every block involved, including all extension-by-zero terms.
Each nonzero summand of the primitive has a join face and a
cut/test block.  Its parent is therefore in both tubes and has
high first-cut length.  On either time endpoint the cut part of
every $\gamma$-face vanishes.  Capping with the primitive and
pushing to the fixed output supplies the required homology
between the two cap orders.

With $\gamma$ first, a $y$-step in its front face gives a
degenerate evaluation.  The surviving shuffles start with
$n+p$ steps from $x\times I$.  Their suffixes are exactly
the shuffles of the retained cut face with $y$.  The native
second cut starts at $x(0)$, so the remaining join cap and
output are precisely those of $\mu(D_Tx,y)$.  Common factor
refinements compare any independently chosen join subdivision.
This proves the first comparison in
\eqref{eq:cj-three-comparisons}.

\paragraph{Symmetry of the joined cut.}
Keep the cut interval whole during the following homotopies.
Rotate the entire loop by $J^{-u}$, obtaining
\[
 (J^{1-u}x)\#(J^{-u}y),\qquad 0\le u\le1.
\]
Then translate the circle parameter of $x\#(J^{-1}y)$
through half a period, ending at $(J^{-1}y)\#x$.
Translation is jointly continuous in $H^1$ and preserves length.
At every stage the evaluations at cut times zero and one are
diagonal.  An intermediate half-time seam need not be diagonal;
any seam terms cancel because we use the full interval.

The swapped join cochain represents the same Thom class relative
to $F_{\eta_j}$, and hence differs from $\sigma$ by the
coboundary of a cochain relative to that far set.  A nonzero
face of this primitive forces its whole adequate parent into
the join tube.  Cap with the primitive followed by the cut and
test cochains and use the boundary formula; the cut kills the
time endpoints.  Shuffle symmetry now identifies the result
with $P(y,x)$.  Repeating the first-half argument in the
opposite order proves the last comparison in
\eqref{eq:cj-three-comparisons}; the second half vanishes
because $N>n+p$.

\paragraph{One length bound for all these homologies.}
Choose the join radius so that each partial rotation changes
the length of every path with the prescribed finite bounds by
less than $\varepsilon$.  This is possible after including
the bounds for any additional finite fillings.  In the whole-loop
homotopies a surviving test saves at least $L-\varepsilon$
from the common joined length, at a cost of two short arcs and
small rotations.  In the first-half homotopy with fixed output,
compare the tested $J^u C_1$ with $C_1$; each differs in length
by less than $\varepsilon$.  Together with the two cut arcs
and the output rotation this gives a cost less than
$8\varepsilon$.  Timing changes preserve length.  The cup-one
and primitive comparisons satisfy the same bound by whole-parent
support and oscillation.  Thus each individual homology lies
already below $a+b-L+1$ after sufficiently small choices.
Their sum, together with the image-preserving refinement and
normalization homologies, satisfies the generous common bound
$a+b-L+4$.  The number of chain terms does not add lengths:
the support of a sum is contained in the union of their supports.
Neither a larger degree nor a finer subdivision traverses an
additional long loop.  This proves the claimed uniformity.
\end{proof}

\subsection{Relative scalar comparison and the absolute unit}

The remaining identification uses the established singular models
of the two string-topology products.  We state the scalar comparison
separately because it does not determine constant-loop output
components.

\begin{lemma}[Comparison with the string-topology products]
\label{lem:cj-public-comparison}
The join $\mu$ induces $\bullet$ on ordinary homology.
If $T$ represents a relative class $[T]\in H^p(\Lambda,M)$
and $b$ is a normalized relative cocycle of degree $r$, then
for every absolute cycle $z$ of degree $p+r+n-1$,
\begin{equation}\label{eq:cj-scalar-comparison}
 \langle b,D_Tz\rangle
   =\langle[T]\circledast[b],[z]\rangle.
\end{equation}
Here the brackets on the right also denote the image of $z$ in
relative homology.
\end{lemma}

\begin{proof}
Hingston--Wahl give singular chain models on $H^1$ loops by
capping with the evaluation pullback of a diagonal Thom cocycle,
then adding short arcs and concatenating or cutting
\cite[Theorem~A, Definitions~1.2 and~1.4]{HingstonWahl2023}.
The dual relative operation is the Goresky--Hingston product
\cite[Theorem~2.13]{HingstonWahl2023}; all sign differences
disappear over $\kk$.  We use these models with the auxiliary
round metric and compare the actual maps and cochains.
Choose their support radius to equal ours and their map domain to
contain our larger tube; both choices can be made within their
injectivity-radius requirement.

For joining, partially rotate the first input by $J^u$ and
apply their short-arc retraction, whose values have equal starts
throughout.
Writing $z=x(0)$ and $z'=y(0)$, the residual mismatch satisfies
$d_{\rm rd}(J^u z,z')=(1-u)d_{\rm rd}(z,z')$,
so all inputs remain in its domain.
At $u=0$ this is their retraction; at $u=1$ the inputs already
have equal starts, so their retraction is the identity and gives
our rotation retraction.  This is a homotopy inside equal-start
pairs.  Interpolate their energy-optimal concatenation times to
halves.  For path pieces of square-root energies $P,Q$, the
optimal time for the first piece is $P/(P+Q)$; during this
interpolation its energy contribution is bounded by
\[
 \max\{2P^2,P(P+Q)\}.
\]
This tends to zero if that piece collapses.  The other piece is
handled in the same way, and if both collapse both contributions
vanish.  Thus the timing homotopy is continuous also at constant
pieces.

For cutting, their short-arc retraction followed by cutting gives
exactly the two paths in \eqref{eq:cj-cut-maps}, with optimal
instead of half timings.  The same interpolation compares them.
At either time endpoint one of the two paths remains constant.
Pairing with $T$ and $b$ evaluates on $z\times I$ the cochain
\[
 U_c(x(0),x(t))\smile\chi,\qquad
 \chi=T(C_1(x,t))\smile b(C_2(x,t)).
\]
Both tests are relative to constants.  Hence $\chi$ vanishes
at $t=0$ through its first factor and at $t=1$ through its
second factor.

Replace $U_c$ by the reduced ordinary Thom representative used
in that model.  They differ by a coboundary relative to the
far set.  Although its primitive need not vanish on the diagonal,
all of its time-end contributions are multiplied by $\chi$
and vanish.  If a cap convention reverses the cup order, the
cup-one formula supplies the comparison.  On a small simplex in
the tube, $\chi$ is a closed pullback; on a simplex in the
far set, every face of the accompanying Thom cochain or relative
primitive vanishes.  Thus extension-by-zero errors contribute
nothing, exactly as in Lemma~\ref{lem:cj-small-chains}.
The same argument permits replacement of either relative test
by a cohomologous representative.  Join Thom representatives
are compared in the same way, without interval-end terms.

Finally the Alexander--Whitney map in the coproduct evaluates
$T\otimes b$ as the displayed cup of the two pullbacks.
Their small-chain replacement and our refined products
represent the same relative chain class: both are homologous
to the input, and small-chain inclusion is a quasi-isomorphism,
also relative to the time ends, by subdivision.  This proves
\eqref{eq:cj-scalar-comparison} and the asserted join
identification.  The argument compares two relative tests;
it makes no assertion about a canonical absolute extension of
the coproduct.
\end{proof}

\begin{proposition}[The cut of the power is the absolute unit]
\label{prop:absolute-unit}
For the choices $S,T,l,L$ in \eqref{eq:cj-thresholds},
\begin{equation}\label{eq:cj-absolute-unit}
 [D_TS]=[M]\quad\text{in }H_n(\Lambda).
\end{equation}
Moreover $D_TS$ is supported in $\mathcal A_{l-L+1}$, and there is a
constant $B$ independent of $m$ such that it is homologous to a
fundamental constant-loop cycle in $\mathcal A_B$.
\end{proposition}

\begin{proof}
Represent $\Theta^{\bullet(m+1)}$ by a finite cycle $z$.
Its cut has degree
$n+2(m+1)d-q=3n-2$.  Lemma~\ref{lem:cj-public-comparison},
commutativity of the relative product over $\kk$, and
\eqref{eq:cj-power-pairing} give
$\langle\theta,D_Tz\rangle=1$.  The rank-one calculation
therefore yields $[D_Tz]=\Theta$.

Both $|z|$ and $|S|=n+q$ exceed $n+p=q+1$.
Apply Theorem~\ref{thm:filtered-interchange} to these cycles,
now forgetting the finite output length bound.  Using the join
identification,
\[
 0\ne\Theta^{\bullet(m+1)}
   =[\mu(D_Tz,S)]=[\mu(D_TS,z)].
\]
If $D_TS$ were null-homologous, its join with $z$ would bound,
by Lemma~\ref{lem:cj-operations}.  Thus $D_TS\ne0$ in
$H_n(\Lambda)=\kk[M]$, proving \eqref{eq:cj-absolute-unit}.
This determines the absolute constant component by interchange
and nonvanishing, rather than by the relative scalar test alone.

The cut bound gives support in $\mathcal A_{l-L+1}$.  Its difference
from a constant fundamental cycle is globally null in degree
$n$; Proposition~\ref{prop:fixed-degree-death} fills that
difference in $\mathcal A_{l-L+1+D_n}$.  By
\eqref{eq:cj-fixed-difference}, this is one fixed level $B$.
\end{proof}

The canonical absolute coproduct of
\cite[Theorem~B]{HingstonWahl2023} removes evaluation components;
the cut used here retains the component determined in
Proposition~\ref{prop:absolute-unit}.  The comparison proved
above is specific to this cut and its degree range.  It is not
an invocation of a general Frobenius identity; the canonical
extension fails such an identity
\cite[Remark~4.14]{HingstonWahl2023}.

We have obtained three compatible pieces of information: a cut
which lowers degree and saves length, an interchange homology
with fixed length cost, and a bounded-length correction of
$D_TS$ to the unit.  The next section combines them with
fixed-degree filling to remove the dependence of the filling
increment on the homological degree.
